\documentclass[11pt,letterpaper]{article}
\usepackage[T1]{fontenc}
\usepackage[utf8]{inputenc}
\usepackage{lmodern,microtype}
\usepackage[margin=1in,headheight=14pt]{geometry}
\usepackage{amsmath,amssymb,amsthm,mathtools}
\usepackage{booktabs,longtable,array,enumitem,fancyhdr,xcolor}
\usepackage{xurl}
\usepackage[colorlinks=true,linkcolor=blue!45!black,urlcolor=blue!45!black,citecolor=blue!45!black]{hyperref}
\usepackage{listings}
\hypersetup{pdftitle={A canonical dual frame for the GRH real-place estimate},pdfauthor={ChatGPT}}
\setlength{\emergencystretch}{2em}
\setcounter{tocdepth}{1}
\setlist{itemsep=3pt,topsep=5pt}
\pagestyle{fancy}
\fancyhf{}
\fancyhead[L]{\small A canonical dual frame for the GRH real-place estimate}
\fancyfoot[R]{\thepage}
\newtheorem{theorem}{Theorem}[section]
\newtheorem{lemma}[theorem]{Lemma}
\newtheorem{proposition}[theorem]{Proposition}
\newtheorem{corollary}[theorem]{Corollary}
\theoremstyle{remark}
\newtheorem{remark}[theorem]{Remark}
\newcommand{\C}{\mathbb C}
\newcommand{\R}{\mathbb R}
\newcommand{\Id}{I}
\newcommand{\tr}{\operatorname{tr}}
\newcommand{\dd}{\,\mathrm d}
\newcommand{\norm}[1]{\left\lVert#1\right\rVert}
\newcommand{\abs}[1]{\left|#1\right|}
\newcommand{\code}[1]{{\small\urlstyle{tt}\nolinkurl{#1}}}
\newcommand{\src}[1]{{\footnotesize\urlstyle{tt}\nolinkurl{#1}}}
\newcommand{\rev}{\code{f2d6ab6c4b70ea73b2deac9e33ad966bc92ea36f}}
\lstset{basicstyle=\ttfamily\footnotesize,breaklines=true,columns=fullflexible,frame=single,rulecolor=\color{black!20},xleftmargin=4pt,xrightmargin=4pt}

\title{\textbf{A canonical dual frame for the GRH real-place estimate}\\[6pt]
  \large Removing triangular normalization while preserving the exact bound}
\author{A technical analysis of the ProveIt Lean development}
\date{November 27, 2026}

\begin{document}
\maketitle
\begin{abstract}
The real-place estimate in the supplied Dirichlet GRH development first chooses an
orthonormal negative frame with triangular pairing matrix, then cancels the
determinants introduced by that normalization against an affine Hua integral.
We give a direct replacement: retain the canonical dual frame, keep its positive
Gram matrix $H$, and integrate a bordered determinant. The identity
\[
 \int_{\C^{k\times i}}
 \det\bigl(H+[X,v]^*[X,v]\bigr)^{-(k+i)}\dd X
 =C_{k,i}\det(H)^{-i}(H_{zz}+\norm v^2)^{-k}
\]
absorbs every normalization factor at once. Here $v$ is fixed,
$C_{k,i}=\prod_{r=0}^{i-1}\pi^k r!/(k+r)!$, and $C_{k,i}\leq1$ when $k\geq7$.
Two short geometric certificates recover exactly the current public bound.
The proposal changes a contained analytic ingredient of the spanning theorem,
with no change to its consumers or numerical hypotheses. It is distinct from the
28 proposal packages present at the inspected commit. The report gives complete
mathematical proofs, source anchors, a migration plan, exact regression checks,
and a separately kernel-checked Lean algebra helper. The full GRH substitution
has not been compiled, so net source and build-time savings remain unmeasured.
\end{abstract}

\tableofcontents
\newpage

\section{Recommendation and scope}

\textbf{Replace the specialized triangular negative-basis layer in
\code{New/Spanning/RemainingTerms/OneRealPlace.lean} by a canonical dual frame and
a bordered Hua integral.} The target is the proof of
\code{New.RealPlace.exists_majorant_volumeRatio_inv_pow}. Its statement can be
retained verbatim~\cite{source}.

The inspected repository is \emph{VladimirReshetnikov/ProveIt}, at commit
\rev. All source paths in this report are relative to
\src{Analysis/GeneralizedRiemannHypothesis/Lean/grh/src/}, unless otherwise stated.
The pinned project uses Lean 4.35.0 and Mathlib revision
\code{56a005855ea17e0b0958f73ec801ae140b0a1679}.

A final comparison with the newer main-branch commit
\code{b47dc5361a5869231f930aa3f59a1c6b53a35f6b} found the entire proposal tree
and the four relevant real-place/Hua/helper files byte-identical. The bundled
\code{artifacts/head_drift.json} records this check. Source counts and references
remain pinned to the earlier revision.

The readable source tree contains 803 Lean files and 476,849 physical lines.
The target module contains 1,030 physical lines, or 712 nonempty lines after
comments are removed. Excluding the unchanged public theorem signature and
documentation, the proposed rewrite touches 437 physical lines / 330
code lines in that module, and makes a 29-line / 26-code-line triangular
normalization helper potentially unnecessary. These are \emph{gross rewrite
targets}, not claimed net deletions. A replacement still needs the canonical
frame lemmas, the general bordered integral, and the final assembly.

The simplification is substantial at the level of the local argument: it removes
an existential basis, its triangularity invariant, its leading principal block,
its last diagonal entry, and the determinant identities reconciling them. It does
not by itself remove a large fraction of the complete GRH development. That
distinction matters in a source tree of this size.

\subsection{What was checked}

The mathematical replacement is proved in Sections~\ref{sec:integral}--\ref{sec:assembly}.
The bundled Lean helper verifies finite-dimensional algebraic identities on the
same Lean/Mathlib versions as the pinned source. It does not formalize the
measure transformation or the entire geometric substitution. The accompanying
Python program passes 717 checks, including Gaussian-rational matrices,
exact rational indefinite-unitary examples, both last-column cases, and a
one-complex-variable integral. These checks supplement the proofs.

This analysis studies the supplied proof and its simplification interfaces. It
does not claim a fresh kernel replay of the full RH/GRH development or a global
novelty result for the underlying matrix-integral identity.

\section{Where the change sits in the proof}

The main comparison uses three inputs: short forms at controlled index, enough
forms to span almost all of the target, and a lower bound on the size of integral
forms. The real-place majorant belongs to the second input. It controls the
remaining terms in the second-moment calculation by a Schur test~\cite{guide,consumer}.

The immediate consumer is
\code{New.Filling.HasCellRank.exists_majorant_lintegral_le_kernelBound}
in \code{SchurOrbit.lean}; its call to the target theorem is at line 553 of the
pinned file. The surrounding argument passes from one real place to all infinite
places, from a lattice sum to an integral, and then to orbit bounds. Consequently
an unchanged local statement and unchanged constant allow the downstream
argument to remain as it stands~\cite{consumer}.

The current local proof takes these steps:

\begin{enumerate}
\item Construct the canonical negative frame $\Omega_0$ from the frame's left
inverse and prescribed pairings.
\item Apply Gram--Schmidt to obtain $\Omega=\Omega_0L$ with
$\Omega^*J\Omega=-\Id$, where $L$ is lower triangular.
\item Express the volume-ratio bound in terms of $\abs{\det L}^2$ and the
positive coordinates of $\Omega'-\Omega$.
\item Split $L$ into its leading block and its last diagonal entry. The
unipotent translation acts through the leading block, so integrating introduces
its determinant as a Jacobian.
\item Prove separate bounds for $\det L$ and $L_{zz}$, and cancel these factors.
\end{enumerate}

The proposed proof keeps Step 1, then works directly with
$H=-\Omega_0^*J\Omega_0$. The last column is fixed by translation without any
triangularity hypothesis. A single integral identity gives the exact powers of
$\det H$ and $H_{zz}+\norm{v_z}^2$ that the geometry controls.

\section{The bordered determinant integral}\label{sec:integral}

Throughout, $*$ denotes conjugate transpose. Complex matrix spaces carry product
Lebesgue measure: one copy of the usual plane measure for each complex entry.
Positive definite Hermitian matrices have positive real determinants, so the
powers below are ordinary real powers.

\begin{theorem}[Bordered Hua identity]\label{thm:border}
Let $k\geq1$, $i\geq0$, and let
\[
 H=\begin{pmatrix}A&b\\b^*&h\end{pmatrix}>0
 \quad\text{in }M_{i+1}(\C).
\]
For fixed $v\in\C^k$, put $c=h+\norm v^2$. Then
\begin{equation}\label{eq:border}
 \int_{\C^{k\times i}}
 \det\bigl(H+[X,v]^*[X,v]\bigr)^{-(k+i)}\dd X
 =C_{k,i}\det(H)^{-i}c^{-k},
\end{equation}
where
\begin{equation}\label{eq:constant}
 C_{k,i}=\pi^{ki}\prod_{r=0}^{i-1}\frac{r!}{(k+r)!}.
\end{equation}
The empty product is 1. In particular, for $k\geq7$ the left-hand side is at most
$\det(H)^{-i}c^{-k}$.
\end{theorem}

\begin{proof}
First suppose $i\geq1$. The Schur complement
\[
 S=A-\frac{bb^*}{h}
\]
is positive definite and $\det H=h\det S$. Set
\[
 U=X-\frac{vb^*}{h},\qquad
 P=\Id_k-\frac{vv^*}{c}.
\]
The eigenvalues of $P$ are 1 on the orthogonal complement of $v$ and $h/c$ on
the line generated by $v$. This also covers $v=0$, and proves $P>0$ and
$\det P=h/c$.

The lower-right entry of $H+[X,v]^*[X,v]$ is $c$. Expanding its Schur complement
gives the exact matrix identity
\begin{align}
 A+X^*X-\frac{(b+X^*v)(b^*+v^*X)}{c}
 &=S+U^*PU.\label{eq:completion}
\end{align}
For example, $Pv=(h/c)v$, and the constant correction on the right is
$-bb^*/h+bb^*\norm v^2/(hc)=-bb^*/c$; the mixed terms agree as well.
Thus
\begin{equation}\label{eq:schurdet}
 \det(H+[X,v]^*[X,v])=c\det(S+U^*PU).
\end{equation}

Make the affine change of variables
$U=P^{-1/2}YS^{1/2}$. Left multiplication by a complex $k\times k$ matrix on
$i$ columns has real Jacobian $\abs{\det P^{-1/2}}^{2i}=\det(P)^{-i}$.
Right multiplication by $S^{1/2}$ on $k$ rows contributes $\det(S)^k$.
The integral therefore equals
\begin{align*}
 &(c\det S)^{-(k+i)}\det(P)^{-i}\det(S)^k
 \int_{\C^{k\times i}}\det(\Id_i+Y^*Y)^{-(k+i)}\dd Y\\
 &\hspace{1cm}=\det(H)^{-i}c^{-k}
 \int_{\C^{k\times i}}\det(\Id_i+Y^*Y)^{-(k+i)}\dd Y.
\end{align*}
The cancellation uses $\det P=h/c$ and $\det S=\det H/h$. The exponent
$k+i$ makes the remaining power of $h$ equal to zero.

For completeness, write
$J_{k,i}(s)=\int\det(\Id_i+Y^*Y)^{-s}\dd Y$.
Sylvester's determinant identity permits replacing $\Id_i+Y^*Y$ by
$\Id_k+YY^*$. Split off the last column $y$ of $Y$, and put
$M=\Id_k+Y_0Y_0^*$. The matrix determinant lemma and $y=M^{1/2}z$ give
\[
 \int_{\C^k}\det(M+yy^*)^{-s}\dd y
 =\det(M)^{-(s-1)}\pi^k\frac{\Gamma(s-k)}{\Gamma(s)},\qquad s>k.
\]
The scalar integral follows directly by writing
$(1+\norm z^2)^{-s}=\Gamma(s)^{-1}
\int_0^\infty t^{s-1}e^{-t}e^{-t\norm z^2}\dd t$ and using Tonelli with the
complex Gaussian integral $(\pi/t)^k$. Iterating over the columns yields
\[
 J_{k,i}(s)=\pi^{ki}\prod_{r=0}^{i-1}
 \frac{\Gamma(s-k-r)}{\Gamma(s-r)},\qquad s>k+i-1.
\]
At $s=k+i$ this is exactly \eqref{eq:constant}. All integrations have
nonnegative integrands, and the displayed gamma values are finite, justifying
both the iteration and convergence. If $i=0$, the integration space is a point,
$\det(H+[v]^*[v])=c$, and \eqref{eq:border} is immediate.

Finally $r!/(k+r)!\leq1/k!$. The quantity $\pi^k/k!$ is below 1 at $k=7$,
and its ratio at successive integers is $\pi/(k+1)<1$. Hence each factor in
\eqref{eq:constant} is at most 1 for $k\geq7$.
\end{proof}

\subsection{The shorter route using the existing Lean Hua bound}

The exact constant is useful for explaining the cancellation, but the integrated
proof need only formalize the inequality. The current development already proves
\code{New.lintegral_norm_det_add_inv_pow_le} in
\code{Lib/General/Analysis/HolomorphicIntegrals.lean}. It handles an affine row
change $A\mapsto RA+a$ and a positive matrix $G$~\cite{hua}.

There is an especially convenient equivalent factorization:
\begin{equation}\label{eq:sylvester}
 \det(H+[X,v]^*[X,v])
 =\det H\;\det\left(\Id_k+\frac{vv^*}{h}+US^{-1}U^*\right).
\end{equation}
Indeed, factoring $H$ by its Schur complement gives
\[
 [X,v]H^{-1}[X,v]^*=US^{-1}U^*+\frac{vv^*}{h},
\]
and Sylvester gives \eqref{eq:sylvester}.
Choose any Gram factor $S^{-1}=R^*R$, and set
$G=\Id_k+vv^*/h$. Then $US^{-1}U^*=(RU^*)^*(RU^*)$.
Such a factor is already available through
\code{CStarAlgebra.isStrictlyPositive_iff_eq_star_mul_self}, used by the
current Hua development. Conjugate transpose is measure preserving, so the
existing affine Hua theorem applies directly to $RU^*$. It yields
\[
 \det(H)^{-(k+i)}\abs{\det R}^{-2k}\det(G)^{-k}
 =\det(H)^{-(k+i)}\det(S)^k(c/h)^{-k}
 =\det(H)^{-i}c^{-k}.
\]
This route reuses the existing integral theorem. It requires neither a new
special-function library nor a triangular basis adapted to the geometric frame.
An arbitrary positive Gram factor inside a generic matrix-integral lemma is
sufficient.

For a direct Lean implementation, state the inequality with row variable
$T\in\C^{i\times k}$ and set $X=X_0+T^*$ from the outset. Then
\[
 RU^*=RT+R\left(X_0^*-\frac{bv^*}{h}\right).
\]
This is literally the affine row form of the existing theorem. Thus the
integration can use the source's existing variable and measure directly;
no separate conjugate-transpose measure-transport lemma is needed.

\begin{remark}[The fixed column and the dimension condition]
Only the first $i$ columns are integrated. Integrating all $i+1$ columns at
exponent $k+i$ reaches the divergent boundary $s=k+(i+1)-1$.
The final column must remain fixed. Also, the unit constant is false in general
for small $k$: $C_{1,1}=\pi$, and $C_{6,1}=\pi^6/6!>1$.
The source's hypothesis $k\geq7$ is exactly adequate for this proof.
\end{remark}

\section{The canonical negative frame}\label{sec:geometry}

We now use the source's real-place setting. Let $J=J_1$ be a Hermitian involution.
Its form has $\varrho$ hyperbolic pairs, a positive part $X_0$ of dimension $k$,
and one negative vector $z$. Let
\[
 F=(e_1,\ldots,e_\varrho,z),\qquad
 F_s=(e_1,\ldots,e_i,z),\qquad n=i+1.
\]
The $e_a$ are isotropic and mutually orthogonal, and $z$ has norm $-1$ and is
orthogonal to them. Let $E_s$ be the column-selection matrix, so $F_s=FE_s$.
Write $E_0$ for the matrix whose columns are the orthonormal coordinate vectors
of $X_0$.

The group $G_1$ consists of the $J$-unitary matrices. For fixed $y\in G_1$, let
$Q_y$ be the source's \code{frameLeftInverse y}. It satisfies $Q_yF=\Id$ and
the columns of $JQ_y^*$ lie in the negative space $yV_-$, where $V_-$ is the
$-1$ eigenspace of $J$.

Fix a matrix of frame pairings
$Z=F^*J\gamma F$, and suppose its selected block
$B=E_s^*ZE_s$ is invertible. Define
\begin{equation}\label{eq:canonical}
 D=ZE_sB^{-1},\qquad
 \Omega_0=JQ_y^*D,\qquad
 H=-\Omega_0^*J\Omega_0.
\end{equation}
For each $\gamma$ with the prescribed $Z$, put
\begin{equation}\label{eq:difference}
 \Omega'_\gamma=\gamma F_sB^{-1},\qquad
 V_\gamma=E_0^*(\Omega'_\gamma-\Omega_0).
\end{equation}
The objects $\Omega_0$ and $H$ depend on $y,Z$, not on the choice of $\gamma$
within its fixed-pairing family.

\begin{lemma}[Canonical pairings and positivity]\label{lem:canonical}
One has
\[
 F^*J\Omega_0=D,\qquad F_s^*J\Omega_0=\Id_n,\qquad H>0.
\]
Moreover
\begin{equation}\label{eq:quotient}
 V_\gamma^*V_\gamma
 =(\Omega'_\gamma-\Omega_0)^*J(\Omega'_\gamma-\Omega_0).
\end{equation}
\end{lemma}

\begin{proof}
The first equality follows from $J^2=\Id$ and $Q_yF=\Id$; the second uses
$E_s^*D=BB^{-1}=\Id$. Thus $\Omega_0$ has independent columns. Put
$W=y^{-1}\Omega_0$. Its columns lie in $V_-$, so
$H=W^*W>0$.

Both $\Omega'_\gamma$ and $\Omega_0$ have pairings $D$ with the full frame.
Their difference is therefore $J$-orthogonal to $F$. A vector orthogonal to all
$e_a$ and to $z$ has zero opposite-isotropic and $z$ coordinates; the remaining
isotropic coordinates contribute nothing to its Gram matrix. Consequently its
$J$-Gram matrix is precisely that of its $X_0$ coordinates, proving
\eqref{eq:quotient}. This is the existing lemma
\code{RealPlace.gram_definiteCols_mul_eq}.
\end{proof}

\subsection{The only retained quotient argument}

Let the source hypothesis \code{NormOnePairOccursFrom j gamma} hold and let
$i=\min(\varrho,j)$. It supplies the shared flag vectors needed to cancel the
unused directions in the determinant of the volume ratio.

The canonical frame is orthogonal to the shared isotropic intersection, as
required by the existing structure \code{IsNegativeBasis}. To see this, if
$Fx=\gamma Fx'$ and $x_z=0$, then invariance of the form gives
$\abs{x'_z}=\abs{x_z}=0$. Hence $Z^*x=0$, and
$D^*x=B^{-*}E_s^*Z^*x=0$.
Together with Lemma~\ref{lem:canonical}, this supplies all three fields of
\code{IsNegativeBasis} without normalizing any vector.

We retain the current quotient-volume lemma
\code{IsNegativeBasis.volumeRatio_eq}; it contains genuine flag cancellation
and should not be deleted. With the selected pairing matrix equal to the
identity, it specializes to
\begin{equation}\label{eq:volume}
 \nu(\gamma,y)=\abs{\det B}\,\abs{\det K},
 \qquad K=\Omega_0^*J\Omega'_\gamma.
\end{equation}
Here $\nu$ is the source's positive \code{volumeRatio}.

\begin{proposition}[Volume bound without normalization]\label{prop:volume}
Under these hypotheses,
\begin{equation}\label{eq:lower}
 \nu(\gamma,y)\ \geq\
 \frac{\abs{\det B}}{2^n}\det(H+V_\gamma^*V_\gamma).
\end{equation}
\end{proposition}

\begin{proof}
The form is negative semidefinite on the span of $\gamma F$, and hence on the
columns of $\Omega'_\gamma$. Expanding \eqref{eq:quotient} gives
\[
 -K-K^*-(H+V_\gamma^*V_\gamma)
 =-(\Omega'_\gamma)^*J\Omega'_\gamma\ \geq\ 0.
\]
Apply the same Ostrowski--Taussky lemma already used in the original proof~\cite{gaussian}:
if $T\geq0$ and $T\leq A+A^*$ for $n\times n$ matrices, then
$\det T\leq2^n\abs{\det A}$. With $T=H+V_\gamma^*V_\gamma$ and $A=-K$, this
gives $\det T\leq2^n\abs{\det K}$. Equation~\eqref{eq:volume} proves the claim.
\end{proof}

\subsection{Translations become ordinary translations}

In the exact family used by the source, the positive coordinates obey
\[
 E_0^*\gamma F_s=E_0^*\gamma_0F_s+T^*B_{<i,*},
 \qquad T\in\C^{i\times k}.
\]
Multiplying by $B^{-1}$ immediately gives
\begin{equation}\label{eq:translation}
 V_\gamma=V_{\gamma_0}+[T^*,0].
\end{equation}
The last column remains fixed, and the other columns have unit-Jacobian affine
coordinates. There is no leading block of $L$ in this formula.

\section{The two geometric certificates}\label{sec:certificates}

\begin{lemma}[Frame-box certificate]\label{lem:box}
If every column of $y^{-1}F$ has Euclidean norm at most $\rho$, then
\begin{equation}\label{eq:detbound}
 \det(H)^{-1}\leq\rho^{2n}.
\end{equation}
\end{lemma}

\begin{proof}
Put $U=y^{-1}F_s$ and $W=y^{-1}\Omega_0$. Since $JW=-W$ and
$F_s^*J\Omega_0=\Id$, we have $U^*W=-\Id$ and $W^*W=H$.
Cauchy--Schwarz for Gram determinants therefore yields
\[
 1=\abs{\det(U^*W)}^2\leq\det(U^*U)\det H.
\]
Hadamard gives $\det(U^*U)\leq\prod_{a=1}^n\norm{U_a}^2\leq\rho^{2n}$.
Divide by $\det H>0$. The needed determinant Cauchy--Schwarz lemma already
appears as \code{New.norm_det_sq_le_of_posSemidef} in the source~\cite{hermitian}.
\end{proof}

Write $v_z$ for the last column of $V_{\gamma_0}$ and set
\[
 c_z=H_{zz}+\norm{v_z}^2.
\]

\begin{lemma}[Lower-floor certificate]\label{lem:lower}
If $j\leq\varrho$, so $i=j$, then $c_z=1$.
\end{lemma}

\begin{proof}
The retained source lemma
\code{exists_common_vector_of_normOnePairOccursFrom} supplies a shared vector
$Fx=\gamma_0F_s a$, with $a_z=1$, $\abs{x_z}=1$, and
$Ba=-x_z e_z$. Therefore
\[
 (\Omega'_{\gamma_0})_z=-\frac{Fx}{x_z}.
\]
This vector lies in the span of $F$ and has $J$-norm $-1$.
Put $d_z=(\Omega'_{\gamma_0})_z-(\Omega_0)_z$. It is orthogonal to $F$, so in
particular orthogonal to $(\Omega'_{\gamma_0})_z$. Equation~\eqref{eq:quotient}
gives $d_z^*Jd_z=\norm{v_z}^2$. Expanding
$(\Omega_0)_z=(\Omega'_{\gamma_0})_z-d_z$ gives
\[
 -H_{zz}=-1+\norm{v_z}^2.
\]
Thus $H_{zz}+\norm{v_z}^2=1$. Notice that this does not assert $H_{zz}=1$.
\end{proof}

\begin{lemma}[Top-floor certificate]\label{lem:top}
If $j=\varrho+1$, so $i=\varrho$, and the source's Bessel length of $y$ is at
most $\theta$, then $c_z\geq\theta^{-2}$.
\end{lemma}

\begin{proof}
At the top the subframe is the whole frame, so $D=\Id$. The vector
$\omega=(\Omega_0)_z$ is orthogonal to every $e_a$ and satisfies $z^*J\omega=1$.
By the existing Bessel-length lemma choose $u$ in the Bessel unipotent subgroup
with $\norm{uy}\leq\theta$. The vector $u^{-1}z$ differs from $z$ by a vector
in the isotropic span of the $e_a$. Hence $(u^{-1}z)^*J\omega=1$.

Move the two vectors by $y^{-1}$. The vector $w=y^{-1}\omega$ lies in $V_-$
and has Euclidean norm $\sqrt{H_{zz}}$. Cauchy--Schwarz gives
\[
 1\leq\norm{y^{-1}u^{-1}z}\sqrt{H_{zz}}
 \leq\theta\sqrt{H_{zz}}.
\]
The second inequality uses the source's norm bound for a frame vector:
for $J$-unitary matrices with $J$ a Hermitian involution, the inverse has the
same operator norm. It follows that $H_{zz}\geq\theta^{-2}$, and adding
$\norm{v_z}^2\geq0$ proves the claim.
\end{proof}

Combining the two cases, for $j\leq\varrho+1$ and $i=\min(\varrho,j)$,
\begin{equation}\label{eq:lastcertificate}
 c_z^{-k}\leq\theta^{2k(j-i)}.
\end{equation}
This replaces the original last-diagonal-entry estimate and the cancellation of
that entry against the rank-one determinant.

\section{Exact recovery of the public theorem}\label{sec:assembly}

Let $s=k+i$, $n=i+1$, $\delta=\abs{\det B}>0$, and
$V_0=V_{\gamma_0}$. Define the explicit majorant
\begin{equation}\label{eq:majorant}
 \phi(T)=\left(\frac{2^n}{\delta}\right)^s
 \det\left(H+(V_0+[T^*,0])^*(V_0+[T^*,0])\right)^{-s}.
\end{equation}
Its determinant is strictly positive for all $T$, so $\phi$ is continuous and
measurable. Equations~\eqref{eq:lower} and \eqref{eq:translation} prove
$\nu(\gamma,y)^{-s}\leq\phi(T)$ whenever $\gamma,T$ meet the source's
translation hypothesis.

Conjugate transpose and translation preserve product Lebesgue measure. Apply
Theorem~\ref{thm:border} to the first $i$ columns of $V_0+[T^*,0]$; its last
column is $v_z$. This proves the stronger identity
\begin{equation}\label{eq:majorantintegral}
 \int\phi(T)\dd T
 =C_{k,i}\left(\frac{2^n}{\delta}\right)^{k+i}
 \det(H)^{-i}c_z^{-k}.
\end{equation}
For $k\geq7$, the two geometric certificates give
\begin{align}
 \int\phi(T)\dd T
 &\leq\frac{2^{(i+1)(k+i)}\rho^{2i(i+1)}\theta^{2k(j-i)}}
 {\delta^{k+i}}\notag\\
 &=\frac{(2\rho^2)^{(i+1)i}
 \bigl(2^{i+1}\theta^{2(j-i)}\bigr)^k}{\delta^{k+i}}.\label{eq:published}
\end{align}
The last expression is exactly the right-hand side of
\code{RealPlace.exists_majorant_volumeRatio_inv_pow}.

\begin{theorem}[Interface-preserving replacement]
Under the hypotheses of the pinned source's
\code{New.RealPlace.exists_majorant_volumeRatio_inv_pow}, the function
\eqref{eq:majorant}, interpreted as an $[0,\infty]$-valued function, satisfies
every conclusion of that theorem, with the same constant. Its construction uses
the canonical frame \eqref{eq:canonical} and does not use
\code{IsTriangularBasis} or \code{exists_isTriangularBasis}.
\end{theorem}

The factor $C_{k,i}$ in \eqref{eq:majorantintegral} also supplies an optional
strict improvement when $i>0$. The recommended first implementation simply
drops it using $C_{k,i}\leq1$, preserving the old contract. We make no claim
that this extra factor alone permits a global change to the rank or scale
hierarchy.

\section{Concrete Lean migration}

\subsection{Retain the geometric reduction that is already useful}

Keep the definitions of \code{frameVolume}, \code{volumeRatio},
\code{frameLeftInverse}, the subframe, and \code{dualProducts}. Keep
\code{IsNegativeBasis}, its quotient-volume theorem, the shared-vector lemma,
\code{norm_det_subframeBlock}, and
\code{frameX0Part_unipotentOf_inv_mul} (whose actual source name uses the
subscript character in $X_0$). They express geometric facts still needed by the
replacement and by \code{SchurOrbit.lean}.

Introduce a canonical definition with the source's existing notation:
\begin{lstlisting}
canonicalNegativeFrame hi y Z :=
  J1 * (frameLeftInverse y).conjTranspose * dualProducts hi Z

canonicalGram hi y Z :=
  -(canonicalNegativeFrame hi y Z).conjTranspose * J1 *
    canonicalNegativeFrame hi y Z
\end{lstlisting}
This is explanatory pseudocode; the delivered compiled helper uses actual Lean
syntax for its generic algebra theorems. An integrated definition should preserve
the source's Unicode identifiers and namespace conventions.

Prove the pairing and positivity lemmas from the first half of the current
\code{exists_isTriangularBasis} proof. That proof already constructs precisely
the matrix $\Omega_0$ in \eqref{eq:canonical}; stop before applying triangular
Gram--Schmidt. Reuse the current proof of orthogonality to the shared isotropic
intersection, now with the right factor $L$ absent.

\subsection{Small reusable lemma list}

A compact replacement API needs these mathematical statements:
\begin{enumerate}
\item Canonical frame pairings and positive Gram matrix.
\item Canonical frame satisfies \code{IsNegativeBasis}.
\item Difference Gram identity and lower bound \eqref{eq:lower}.
\item Translation rule \eqref{eq:translation}.
\item Determinant bound \eqref{eq:detbound}.
\item Lower-floor equality $c_z=1$ and top-floor inequality $c_z\geq\theta^{-2}$.
\item Generic bordered Hua inequality, derived from the existing affine Hua theorem.
\item Assembly \eqref{eq:majorant}--\eqref{eq:published} with the original theorem signature.
\end{enumerate}

The statement of the generic integral should quantify over an arbitrary positive
Hermitian $(i+1)\times(i+1)$ matrix, a fixed last column, and product complex
Lebesgue measure. It should not mention \code{Shape}, $G_1$, Bessel subgroups,
or the GRH parameters. This separation is what makes the simplification reusable.

\subsection{The exact source boundary}

\begin{center}\small
\begin{tabular}{p{2.5in}rrp{1.45in}}
\toprule
Pinned source region & Physical & Code & Treatment\\
\midrule
\code{OneRealPlace.lean}, 449--602 &154&106&Triangular existence and estimates\\
\code{OneRealPlace.lean}, 604--705 &102&70&Normalized image and volume proof\\
\code{OneRealPlace.lean}, 802--916 &115&90&Last-column estimates\\
\code{OneRealPlace.lean}, 941--1006 &66&64&Majorant proof body only\\
\midrule
Subtotal &437&330&Rewrite, not net deletion\\
\code{HermitianMatrices.lean}, 71--99 &29&26&Potentially retire unused helper\\
\bottomrule
\end{tabular}
\end{center}

These totals count the public theorem's proof body without counting its retained
statement. Together with the optional helper, the gross target is 466 physical
lines / 356 code lines. The complete source audit scans all 803 Lean files and
records no outside lexical uses of
the specialized triangular frame, its existence lemma, \code{imageBasis}, or
the normalized difference construction. The generic triangular normalization
helper has one source call, inside the target construction. It may be removed
if the new generic integral uses a nontriangular positive Gram factor.

The source boundary is deliberately narrow. In particular, the orthonormal bases
in \code{New/Spanning/Frames/RealPlaceGauge.lean} serve a different measurable
kernel-gauge construction; their removal does not follow from this argument.
The general Hermitian-matrix, Gaussian, Hua, and bordered-link libraries remain
dependencies. No deletion of a whole library is claimed.

\subsection{Integration and acceptance order}

First prove the generic bordered inequality with explicit measure conventions.
Then introduce the canonical frame and replace the local proof while retaining
the old public statement. Rebuild the target and its direct consumer
\code{SchurOrbit.lean}. Finally rebuild the main theorem dependency and compare
the printed theorem types and axiom sets. Only after successful integration
should obsolete declarations be removed and net line counts or build times
reported. The standalone helper's successful compilation is not a substitute
for those integration steps.

\section{Why the proposal is different from the existing ones}

The comparison corpus at \rev{} contains 28 directories, 43 TeX files totaling
11,944 lines, and 136 TeX/Markdown/text files. The audit reviews package summaries,
mathematical mechanisms, targeted integration passages, and source targets. It
does not claim to have independently re-proved every earlier proposal. A hashed
text-file inventory and a complete mechanism matrix are bundled~\cite{comparison}.

The nearest apparent neighbors are:
\begin{itemize}
\item \code{gram_adjoint_simplification}: integral adjoints and Weyl--Fock Gram
identities in a one-lift arithmetic height argument. It removes a different
tight-frame construction; it does not evaluate this real-place integral with
the full negative-frame Gram matrix.
\item \code{grh_simplification}: scalar-height/chain architecture and branching
representation norms. Its Gram identities concern extended tensors, not the
fixed-last-column majorant.
\item \code{projective_incidence_grh}: a finite-field positive-type singular-minor
estimate used in a second-moment bound. The present change is in the
archimedean remaining-term integral.
\item \code{centered_theta_lift_simplification}: a centered covariance and a better
rank/completion budget. It changes how the moment estimate is used; the present
proposal changes a local analytic ingredient supplying that estimate.
\end{itemize}

The mechanisms can coexist. In particular, improving the final rank budget does
not make the canonical-frame identity a duplicate, and the proposed replacement
does not count savings already attributed to an earlier package.

\section{Verification record and practical limitations}

\subsection{Mathematical and computational checks}

The report proves the general integral, the canonical-frame volume inequality,
the determinant certificate, both last-column cases, the translation rule,
and exact recovery of the original constant. A second mathematical review
independently derived the same Schur complement, Jacobian powers, and geometric
certificates.

The delivered test program exercises genuinely complex rational entries,
including off-diagonal coupling $b\ne0$; it does not test only diagonal Gram
matrices. Its geometric tests produce rational $J$-unitary transformations by
a Cayley transform and compute the canonical frame directly. They verify the
orthogonality and Gram identities, the reduced volume determinant, and the
lower-floor equality. The top-floor regression uses a weaker Frobenius-norm
certificate for the identity Bessel element; the full Bessel-length statement
is supplied by the proof, not inferred from those examples.

The 717 checks comprise 147 block/Schur/Jacobian checks, eight Cayley-unitarity
checks, 41 canonical-geometry and boundary checks, 504 rational upper bounds
for Hua constants, one radial numerical integral, and 16 exact checks of the
published constant. The radial example uses $k=i=1$, $h=3$, $v=2$, $A=2$, $b=1$;
its integral is $\pi/35$. It deliberately lies outside the unit-constant
$k\geq7$ range, testing the exact formula rather than a vacuous bound.

\subsection{Formal status}

\code{formal/BorderedGram.lean} is a standalone checked component.
Its accompanying verification log and README identify the exact proved
declarations and pinned dependencies. The component establishes the algebraic
completion and determinant/cancellation identities used here. It adds no
assumption asserting RH, GRH, or the desired integral inequality. All six
theorems compile without warnings; each printed axiom set contains only
\code{propext}, \code{Classical.choice}, and \code{Quot.sound}. In particular,
none depends on \code{sorryAx}.

The following parts remain to be connected in the GRH project: positive
definiteness and factorization in the generic bordered integral, its block
determinant rewrite and application of the existing affine Hua bound, the
canonical geometric lemmas in the source's types, and replacement of the
original majorant proof. The row-variable formulation above uses the source's
existing product measure directly. Their mathematical
arguments are given above. They have not been represented as already compiled
Lean theorems.

The recommended next implementation is therefore concrete and reviewable:
one generic integral lemma and one local normalization refactor, preserving
\eqref{eq:published}. The archive supplies the proof of that replacement and
the evidence needed to assess it, without presenting anticipated net savings
as measured results.

\appendix
\section{Pinned source references}

The reference links below resolve to the inspected commit. The artifact
\code{artifacts/source_manifest.json} records hashes, while
\code{code/verify_sources.py} checks those bytes in an available Git checkout.

{\small
\begin{thebibliography}{9}
\setlength{\itemsep}{3pt}
\bibitem{source}
\href{https://github.com/VladimirReshetnikov/ProveIt/blob/f2d6ab6c4b70ea73b2deac9e33ad966bc92ea36f/Analysis/GeneralizedRiemannHypothesis/Lean/grh/src/New/Spanning/RemainingTerms/OneRealPlace.lean}{OneRealPlace.lean: canonical precursor, triangular basis, volume bound, and public majorant.}

\bibitem{hua}
\href{https://github.com/VladimirReshetnikov/ProveIt/blob/f2d6ab6c4b70ea73b2deac9e33ad966bc92ea36f/Analysis/GeneralizedRiemannHypothesis/Lean/grh/src/Lib/General/Analysis/HolomorphicIntegrals.lean}{HolomorphicIntegrals.lean: Hua integral and affine matrix integral bound.}

\bibitem{hermitian}
\href{https://github.com/VladimirReshetnikov/ProveIt/blob/f2d6ab6c4b70ea73b2deac9e33ad966bc92ea36f/Analysis/GeneralizedRiemannHypothesis/Lean/grh/src/Lib/General/LinearAlgebra/HermitianMatrices.lean}{HermitianMatrices.lean: determinant Cauchy--Schwarz and triangular Gram factor.}

\bibitem{gaussian}
\href{https://github.com/VladimirReshetnikov/ProveIt/blob/f2d6ab6c4b70ea73b2deac9e33ad966bc92ea36f/Analysis/GeneralizedRiemannHypothesis/Lean/grh/src/Lib/General/Analysis/GaussianIntegral.lean}{GaussianIntegral.lean: Ostrowski--Taussky determinant inequality.}

\bibitem{consumer}
\href{https://github.com/VladimirReshetnikov/ProveIt/blob/f2d6ab6c4b70ea73b2deac9e33ad966bc92ea36f/Analysis/GeneralizedRiemannHypothesis/Lean/grh/src/New/Spanning/RemainingTerms/SchurOrbit.lean}{SchurOrbit.lean: direct consumer and passage to orbit estimates.}

\bibitem{guide}
\href{https://github.com/VladimirReshetnikov/ProveIt/blob/f2d6ab6c4b70ea73b2deac9e33ad966bc92ea36f/Analysis/GeneralizedRiemannHypothesis/docs/guide/theta_lifts_and_dirichlet_zeros.tex}{Human-readable guide, especially the spanning and volume-comparison sections.}

\bibitem{comparison}
\href{https://github.com/VladimirReshetnikov/ProveIt/tree/f2d6ab6c4b70ea73b2deac9e33ad966bc92ea36f/Analysis/GeneralizedRiemannHypothesis/docs/proposals}{Proposal corpus: all 28 packages in the nonduplication audit.}
\end{thebibliography}
}

\end{document}
